En este capítulo se detalla la situación actual que origina el desarrollo o la mejora de un sistema informático. Asimismo, se describen las alternativas tecnológicas existentes.

\section{Contexto} 
Los trabajos de desarrollo de software pretenden dar una solución a la problemática existente en una determinada organización o de un conjunto de personas. Por ello es preciso describir la situación actual, incluyendo información sobre el entorno tecnológico existente (si lo hubiera). En tal caso, es preciso analizar las fortalezas y debilidades del mismo.

Adicionalmente, puede incluirse información sobre los modelos de procesos de negocio actuales existentes en la organización. Un proceso de negocio no es más que es el conjunto de actividades o tareas relacionadas entre sí y estructuradas en una secuencia específica para producir un servicio o producto (cumple un determinado objetivo de negocio) para un cliente en particular. Pueden emplearse notaciones como BPMN o diagramas de actividades UML.

\section{Estado de la técnica} 
Antes de comenzar con la planificación del proyecto, es preciso realizar un estudio de las invenciones existentes que permitan solventar la problemática descrita en la sección anterior. Para ello se utilizarán buscadores web generales, como Google; buscadores especializados en trabajos científicos, como Google Scholar; y repositorios de aplicaciones, como Apple App Store.

Este estudio permitirá, por un lado, descubrir el grado de novedad de la nueva solución informática a desarrollar y, por otro lado, descubrir soluciones ya existentes (habitualmente de software libre) que puedan servir de plataforma base para el nuevo desarrollo.




